\documentclass[10pt,a4paper]{amsart}
\usepackage[margin=19mm]{geometry}
\usepackage{amsmath,amssymb,mathtools}
\usepackage[scr=rsfs,cal=euler]{mathalfa}
\usepackage{xfrac}
\usepackage[unicode,hidelinks,bookmarks=false]{hyperref}
\newcommand{\ie}{i.e.\ }
\newcommand{\cf}{cf.\ }
\newcommand{\resp}{respectively\ }
\newcommand{\Subsection}{\subsection}
\providecommand{\nobreakdash}{\nobreak\hbox{-}}
\setlength{\emergencystretch}{2em}
\multlinegap0pt
\usepackage{amssymb}
\usepackage{mathtools}
\newtheorem{proposition}{Proposition}
\newcommand{\CS}{\operatorname{CS}}
\newcommand{\im}{\operatorname{im}}
\title{Chern--Simons Fluctuations and Information Geometry\\ in Discrete Electromagnetism}
\author{Jean-Pierre Magnot}
\date{Source: arXiv:2609.00020v1 (CC BY 4.0)}
\begin{document}
\maketitle

\noindent\emph{Source and licence.} Chern--Simons Fluctuations and Information Geometry\\ in Discrete Electromagnetism. Version arXiv:2609.00020v1 (CC BY 4.0), distributed by arXiv under the Creative Commons Attribution 4.0 International licence, http://creativecommons.org/licenses/by/4.0/, as recorded in the arXiv licence metadata for this exact version and shown on the abstract page https://arxiv.org/abs/2609.00020v1. Full licence notice: \texttt{licenses/arxiv-2609.00020-CC-BY-4.0.txt}.\par\smallskip

\noindent\textit{Editorial scope.} Counted unit: the article's proposition prop:gauge-harmonic (source0.tex lines 445--472). The article's complete original proof of it is delivered with it (source0.tex lines 474--556, one \texttt{proof} environment of 83 lines). No mathematics is written, repaired or re-derived: the statement and the proof are the article's own lines, in the article's own notation, and no retained line is substituted or re-flowed.  Retained uncounted as the source-local context the proof needs: lines 253--258; lines 321--330.  Nothing was omitted from any retained span, so the package carries no omission record.  No retained line contains a citation, so the wrapper carries no bibliography and no bibliography entry.  No retained line contains a figure environment, an image-inclusion command or drawing code, so no image is delivered and the package declares no asset dependency.  Editorial formatting only, in addition: an AMS \texttt{amsart} preamble that restates the article's own declarations and macros used by the retained lines, namely the environments \texttt{proposition} on separate counters, together with the standard packages those lines need.  The retained environments print in this document as Proposition 1 in order of appearance, whereas the article prints the same items with the numbering of its own full document.  The complete native source of the article as distributed in the arXiv e-print for this exact version is bundled unchanged as \texttt{sources/208993/source0.tex}.
\begin{equation}
dW_k
=
W_{k+1}\delta_k.
\label{eq:whitney-commutation}
\end{equation}

\begin{equation}
C^1(K)
=
\im\delta_0
\oplus
\mathcal H_K^1
\oplus
\im\delta_1^\ast,
\label{eq:hodge}
\end{equation}

\begin{proposition}
\label{prop:gauge-harmonic}
The bilinear form \(c_K\) is symmetric and satisfies
\begin{equation}
c_K(\delta_0\chi,a)=0
\label{eq:exact-annihilation}
\end{equation}
for all \(\chi\in C^0(K)\) and \(a\in C^1(K)\).

It also satisfies
\begin{equation}
c_K(h,a)=0
\label{eq:harmonic-annihilation}
\end{equation}
for every \(h\in\mathcal H_K^1\).

Consequently, if
\begin{equation}
a=a_{\mathrm g}+a_{\mathrm h}+a_{\mathrm T}
\end{equation}
is the decomposition \eqref{eq:hodge}, then
\begin{equation}
\CS_K(a)
=
\CS_K(a_{\mathrm T}).
\label{eq:only-transverse}
\end{equation}
\end{proposition}

\begin{proof}
Let
\begin{equation}
\alpha=W_1(a),
\qquad
\eta=W_1(b).
\end{equation}
Since \(M\) has no boundary,
\begin{equation}
0
=
\int_M d(\alpha\wedge\eta)
=
\int_M d\alpha\wedge\eta
-
\int_M\alpha\wedge d\eta.
\end{equation}
Because forms of degrees two and one commute under the wedge
product,
\begin{equation}
d\alpha\wedge\eta
=
\eta\wedge d\alpha.
\end{equation}
Therefore,
\begin{equation}
\int_M\alpha\wedge d\eta
=
\int_M\eta\wedge d\alpha,
\end{equation}
which proves symmetry.

From \eqref{eq:whitney-commutation},
\begin{equation}
W_1(\delta_0\chi)
=
dW_0(\chi).
\end{equation}
Hence
\begin{align}
c_K(\delta_0\chi,a)
&=
\int_M
dW_0(\chi)\wedge dW_1(a)
\nonumber\\
&=
\int_M
d\left(
W_0(\chi)dW_1(a)
\right)
\nonumber\\
&=
0.
\end{align}

If \(h\in\mathcal H_K^1\), then
\begin{equation}
\delta_1h=0,
\end{equation}
and therefore
\begin{equation}
dW_1(h)
=
W_2(\delta_1h)
=
0.
\end{equation}
Using symmetry,
\begin{equation}
c_K(h,a)
=
c_K(a,h)
=
\int_M
W_1(a)\wedge dW_1(h)
=
0.
\end{equation}

Expanding \(c_K(a,a)\) along \eqref{eq:hodge}, all terms
containing \(a_{\mathrm g}\) or \(a_{\mathrm h}\) vanish.
This gives \eqref{eq:only-transverse}.
\end{proof}

\end{document}
